% Intended LaTeX compiler: lualatex
\documentclass[version=last,draft=true,paper=A4,portrait,twoside=true,twocolumn=false,headinclude=false,footinclude=false,fontsize=11,BCOR=15mm,DIV=calc,pagesize=auto,titlepage=firstiscover,mpinclude=true,open=right,chapterprefix=true,numbers=autoendperiod,headsepline=false,parskip=false]{scrbook}
\usepackage{graphicx}
\usepackage{grffile}
\usepackage{longtable}
\usepackage{wrapfig}
\usepackage{rotating}
\usepackage[normalem]{ulem}
\usepackage{amsmath}
\usepackage{textcomp}
\usepackage{amssymb}
\usepackage{capt-of}
\usepackage{hyperref}
\usepackage{polyglossia}
\setmainlanguage{english}
\setotherlanguages{french,italian}
\usepackage[autostyle=true,english=american,french=guillemets,thresholdtype=words,threshold=3]{csquotes}
\MakeOuterQuote{"}
\MakeForeignQuote{french}{«}{»}
\AtBeginEnvironment{quote}{\itshape}
\usepackage{amsmath}
\usepackage{amsthm}
\usepackage{amssymb}
\usepackage{centernot}
\usepackage{hyperref}
\hypersetup{colorlinks,urlcolor=bibleblue,linkcolor=bibleblue,citecolor=bibleblue,filecolor=black}
\usepackage{balance}
\usepackage{array}
\usepackage{tabularx}
\usepackage{booktabs}
\usepackage[most]{tcolorbox}
\usepackage[french]{fmtcount}
\fmtcountsetoptions{french=france}
\usepackage[singlespacing]{setspace}
\usepackage[super]{nth}
\usepackage[protrusion=true,tracking=true]{microtype}
\usepackage{ragged2e}
\usepackage[all]{nowidow}
\usepackage{enumitem}
\usepackage{adforn}
\usepackage{float}
\usepackage{xcolor}
\usepackage{graphicx}
\graphicspath{ {/home/sync0/Dropbox/paris_1/} }
\usepackage{lipsum}
\usepackage{tikz}
\usetikzlibrary{positioning}
\tikzset{main node/.style={circle,fill=gray!45,draw,minimum size=0.5cm,inner sep=0pt},}
\usepackage{fontspec}
\usepackage{unicode-math}
\usepackage[oldstyle]{libertine}
\defaultfontfeatures{Scale=MatchLowercase}
\newfontfamily\titlefamily[Scale=2]{Linux Biolinum O}
\newcommand\HUGE{\fontsize{30}{33}\selectfont}
\usepackage{scrlayer-scrpage}
\pagestyle{scrheadings}
\clearscrheadfoot
\automark[chapter]{part}
\cehead{vers la machine à gourverner}
\cohead{\MakeLowercase{\headmark}}
\ohead{\pagemark}
\AfterTOCHead{\singlespacing}
\setkomafont{labelinglabel}{\normalfont\sffamily\bfseries}
\setkomafont{minisec}{\usekomafont{subsection}}
\setkomafont{pagehead}{\normalfont\sffamily\mdseries\scshape}
\setkomafont{pagenumber}{\large\rmfamily\upshape}
\setkomafont{sectioning}{\sffamily\mdseries}
\setkomafont{caption}{\small}
\setkomafont{captionlabel}{\sffamily\mdseries\scshape\lowercase}
\setkomafont{part}{\HUGE\scshape\lowercase}
\renewcommand*{\partformat}{}
\renewcommand\partmarkformat{}
\setkomafont{chapter}{\Huge\scshape\lowercase}
\renewcommand{\raggedchapter}{\centering}
\renewcommand*\chapterformat{\thechapter\autodot\par\enskip}
\renewcommand*\chaptermarkformat{}
\setkomafont{section}{\Large\scshape\lowercase}
\setkomafont{subsection}{\large}
\renewcommand*{\subsubsectionformat}{\sffamily\mdseries\upshape \thesubsubsection\autodot\enskip}
\setkomafont{subsubsection}{\large\itshape}
\AtBeginDocument{\renewcaptionname{english}\contentsname{contents}}
\addtokomafont{chapterentry}{\large\sffamily\mdseries\scshape\lowercase}
\usepackage{tocstyle}
\settocfeature{raggedhook}{\raggedright}
\selecttocstyleoption{tocgraduated}
\usetocstyle{nopagecolumn}
\newtcolorbox{note}[2][]{grow to right by=0mm,grow to left by=-1em,boxrule=0pt,boxsep=0pt,opacityback=0.0,breakable,parbox=false,enhanced jigsaw,borderline west={4pt}{0pt}{lightgrey},title={#2},coltitle={black},fonttitle={\sffamily},attach title to upper={},halign title=right,after title={\smallskip\par}#1}
\newtcolorbox{question}[2][]{grow to right by=0mm,grow to left by=-1em,boxrule=0pt,boxsep=0pt,opacityback=0.0,breakable,parbox=false,enhanced jigsaw,borderline west={4pt}{0pt}{darkgrey},title={#2},coltitle={black},fonttitle={\sffamily},attach title to upper={},halign title=right,after title={\smallskip\par}#1}
\newtcolorbox{definition}[3][]{grow to right by=0mm,grow to left by=-1em,boxrule=0pt,boxsep=0pt,opacityback=0.0,breakable,enhanced jigsaw,borderline west={4pt}{0pt}{midgrey},title={#2},coltitle={black},fonttitle={\sffamily\bfseries},fontupper={\normalfont},fontlower={\itshape},lower separated=false,attach title to upper={},after title={\hspace{1em}{\rmfamily\mdseries\itshape #3}\par}#1}
\renewcommand*\labelitemi{\adforn{33}}
\renewcommand*\labelitemii{\adforn{73}}
\renewcommand*\labelitemiii{\adforn{73}}
\renewcommand*\labelitemiv{\adforn{73}}
\definecolor{bibleblue}{HTML}{00339a}
\definecolor{whitegrey}{HTML}{f7f7f7}
\definecolor{lightgrey}{HTML}{cccccc}
\definecolor{midgrey}{HTML}{969696}
\definecolor{darkgrey}{HTML}{636363}
\definecolor{blackgrey}{HTML}{252525}
\newcommand{\notimplies}{\centernot\implies}
\setcounter{secnumdepth}{3}
\setcounter{tocdepth}{1}
\recalctypearea
\author{Carlos Alberto Rivera Carreño}
\date{}
\title{}
\hypersetup{
 pdfauthor={Carlos Alberto Rivera Carreño},
 pdftitle={},
 pdfkeywords={},
 pdfsubject={},
 pdfcreator={Emacs 26.1 (Org mode 9.2)}, 
 pdflang={English}}
\begin{document}

\begin{titlepage}
 \centering
\begin{french}
 {\large \textsc{université paris i panthéon sorbonne} \par}
 \vspace{0.5\baselineskip}
 {\large \textsc{ufr} 02 : Sciences économiques  \par}
 \vspace{0.5\baselineskip}
 {\large Master 2 : Économie et sciences humaines \par}
 \vspace{0.5\baselineskip}
 {\large 2018-2019 \par}
\end{french}
 \vspace{6\baselineskip}
 {\huge\sffamily \textsc{vers la machine à gouverner}  \par}
 {\Large\sffamily Herbert Simon and the Impossibility of a Democratic Computer \par}
\vfill
\begin{french}
 {\large Présenté et sountenu par : \par}
\end{french}
 {\Large Carlos Alberto Rivera Carreño \par}
 \vspace{1\baselineskip}
\begin{french}
 {\large Directeur de mémoire : \par}
\end{french}
 {\Large Jean-Sébastien Lenfant \par}
\end{titlepage}

\pagestyle{empty}

\begin{french}
L'Université Paris 1 Panthéon Sorbonne n'entend donner aucune approbation,
ni désapprobation aux opinions émises dans ce mémoire ; elle doivent être
considérées comme propres à leur auteur. 
\end{french}
\vfill

\newpage
\vspace*{\fill}
\noindent
\includegraphics[height=1.5cm]{gpl3.png}\par
\vspace{1\baselineskip}
This text is free: you can redistribute it and/or modify it
under the terms of the \textsc{gnu} General Public License as published by
the Free Software Foundation, either version 3 of the License or any later
version.

This text is distributed in the hope that it will be useful, but \textbf\{without
any warranty\}; without even the implied warranty of \textbf\{merchantability or 
fitness for a particular purpose\}. See the \textsc{gnu} General 
Public License for more details.

You should have received a copy of the \textsc{gnu} General Public License along
with this text. If not, see \url{http://www.gnu.org/licenses/}.

\vspace{1\baselineskip}
\noindent
Copyright \textcopyright \textsc{sync0} 2018. 

\newpage 
\begin{FlushRight}
\begin{italian}
\textit{Per il professore Giorgio Israel. \newline Benché i nostri destini fossero uniti da quell'anno fatale del 1492, \linebreak ho mancato il nostro incontro. \linebreak Sono arrivato in ritardo, come la coscienza della nostra generazione. \linebreak I campi sanguinanti sono così prossimi, \linebreak ma nessuno vede niente. \linebreak Nell'era dell'informazione, diventiamo più ignoranti. \linebreak Questo è il prezzo del biglietto d'ingresso. \linebreak Grazie a Lei l'ho capito, maestro.}
\end{italian}
\end{FlushRight}

\newpage
\tableofcontents 

\frontmatter
\pagestyle{plain}
\chapter{Note aux lecteurs français} 
\lipsum
\chapter{Acknowledgements} 
\lipsum

\chapter{Introduction} 

A story that although was present for a long time, seems today relegated to
the confines of a few books on the history of computing. 

De Prony's Tables

Simon and Babbage

Simon and Industrial Democracy

\lipsum
\mainmatter
\pagestyle{scrheadings}
\chapter{Introduction}
\label{sec:org38f9efc}
\section{De Prony's Tables}
\label{sec:org4c3c909}
\section{Simon and Babbage}
\label{sec:org0413090}
\section{Simon and Industrial Democracy}
\label{sec:orgc8721e0}
\chapter{A History of the Concept of Labor}
\label{sec:orgfed6883}
\section{From Political Economy to Mechanics}
\label{sec:org0db3e41}
\lipsum
\chapter{Herbert Simon and the Thinking Machines}
\label{sec:orgd0e89ee}
\section{The Cold War and the Computer}
\label{sec:org6e81264}
\lipsum
\section{Artificial Intelligence}
\label{sec:org217d965}
\lipsum
\section{Administrative Behavior}
\label{sec:org19668d2}
\lipsum

\backmatter
\chapter{Conclusion} 
\lipsum
\chapter{References} 
\lipsum
\end{document}